\documentclass[a4paper,12pt]{article}
\usepackage[spanish,activeacute]{babel}
\usepackage[latin1]{inputenc}

%% Para las columnas
\usepackage{multicol}
\usepackage{amssymb,amsmath,eepic,fancyhdr}
\usepackage{latexsym}
\usepackage{datenumber}


%% Para numeraci?n autom?tica
\newcounter{nroproblema}
\setcounter{nroproblema}{1}
\newcounter{nroapartado}
\setcounter{nroapartado}{1}
\newcommand{\n}{%
    \vspace{.1in}\noindent{\bf \arabic{nroproblema}.\ }%
    \addtocounter{nroproblema}{1}%
    \setcounter{nroapartado}{1}%
    }%
\newcommand{\apartado}{\vspace*{0.2cm}\hspace*{0.2cm}%
   {\bf \alph{nroapartado})}
   \addtocounter{nroapartado}{1}%
   }

% \newcommand{\A}{{\bf (A)\; }} % De problemas de Andrei
\newcommand{\A}{} % De problemas de Andrei

\newcommand{\R}{{\mathbb{R}}}
\newcommand{\C}{{\mathbb{C}}}
\newcommand{\N}{{\mathbb{N}}}
\newcommand{\K}{{\mathbb{K}}}


\newcommand{\rpol}[1]{\mathcal P^{\R}_{#1}}
\newcommand{\kpol}[1]{\mathcal P^{\K}_{#1}}

\newcommand{\cuatrox}[1][x]{\mbox{$\Bigg(\begin{array}{c}
    #1_1\\[-4pt] #1_2\\[-4pt] #1_3\\[-4pt] #1_4\end{array}\Bigg)$}}
\newcommand{\cuatroy}{\cuatrox[y]}
\newcommand{\cuatroz}{\cuatrox[z]}

% \newcommand{\cpol}[1]{\mbox{$\mathcal P^{\C}_{#1}}$}




%% m?genes

\setlength{\unitlength}{1mm} \addtolength{\voffset}{-18mm}
\addtolength{\textheight}{35mm} \addtolength{\hoffset}{-18mm}
\addtolength{\textwidth}{40mm}





\begin{document}


\begin{flushright}
1${}^o$ de grado en F�sica, curso 2013/14

\today
\end{flushright}


\begin{center}{\bf {\Large �LGEBRA I. HOJA 2}}\end{center}

% \vskip.5cm


%\noindent{Nota: Las hojas est\'an colgatas tambi\'en en la p\'agina web de del curso
%\begin{verbatim}http://www.uam.es/personal_pdi/ciencias/mzambon/WS12.html\end{verbatim}}



\n \A
Estudiar si los siguientes conjuntos son subespacios vectoriales
de $\R^4$.
\begin{enumerate}\itemsep=-2pt
\item $A=\{(3x_2,x_2,x_4+x_2,x_4)^T\;|\;x_2,x_4\in \R\}$.

\item $B=\{(x_1,x_2,x_3,x_4)^T\;|\;x_1\cdot x_2=0\}$.

\item $C=\{(x_1,x_2,x_1,x_2)^T\;|\;x_1=1\}$.

\item $D=\{(x_1,x_1,x_1,x_1)^T\;|\;x_1\in \R\}$.

\item $E=\{(x_1,x_2,x_3,x_4)^T\;|\; x_1=t,x_2=t,x_3=2t,x_4=5t,
t\in\R\}$.
\end{enumerate}


\n \A
Comprueba si son o no espacios
vectoriales sobre $\mathbb{R}$ los siguientes conjuntos con las operaciones
indicadas:
\begin{enumerate}
\itemsep=-2pt
\item $V=\{(x_1,x_2)^T\in \R^2\},\quad
(x_1,x_2)^T+(y_1,y_2)^T=(x_1+y_1,0),\, \lambda(x_1,x_2)=(\lambda
x_1,0)$.

\item $V=\{(x_1,x_2)^T\in \R^2\mid x_1\not=0, x_2\not=0\},\quad
(x_1,x_2)^T+(y_1,y_2)^T=(x_1y_1,x_2y_2)^T,\, \newline \lambda(x_1,x_2)^T=(\lambda
x_1,\lambda x_2)^T$.
\end{enumerate}

\n
Sea $\K$ el cuerpo de los n\'umeros reales � bien el cuerpo de los n\'umeros complejos
(es decir, $\K=\R$ \'o $\K=\C$).
Dado un $n\in \N$,
denotamos por $\kpol{n}$ el conjunto de polinomios de grado $\le
n$ con coeficientes en $\K$.

Decidir razonadamente si los siguientes conjuntos, con la adici\'on
y multiplicaci�n por escalares en $\K$ definidas de manera natural,
forman espacios vectoriales:
\begin{enumerate}
\item  $\kpol{n}$.
\item  el conjunto de polinomios en $\kpol{n}$ de grado \emph{igual} a $n$.
\item   $\{f\colon [0,1]\to \mathbb{R} \text{ tal que } f \text{ es continua}\}$ ($\K=\R$).
\item   $\{f\colon [0,1]\to \mathbb{R}_+ \text{ tal que } f \text{ es continua}\}$
($\K=\R$), donde $\mathbb{R}_+$ denota los n\'umeros positivos.
\end{enumerate}

\n Sea ${\mathcal P}^\R$ el espacio de todos los polinomios  sobre $\mathbb{R}$. Comprobar que, con las  adici\'on y multiplicaci\'on por escalares
definidas de manera natural,  es un espacio vectorial.
?`Existe una base de ${\mathcal P}^\R$?


\n �Cu�les de las siguientes familias de vectores forman un sistema de generadores de $\mathbb{R}^3$
(con las operaciones usuales)?
\begin{enumerate}
\itemsep=-2pt
\item $\{(1,2,3)^T,(1,2,5)^T\}$,
\item $\{(1,2,3)^T,(1,2,5)^T,(0,0,-4)^T\}$,
\item $\{(1,2,3)^T,(1,2,5)^T,(0,0,-4)^T,(1,1,1)^T\}$.
\end{enumerate}

\n �Cu�les de las 3 familias de vectores que aparecen en el ejercicio anterior son linealmente independientes?

\n �Cu�les de las 3 familias de vectores que aparecen en el ejercicio anterior forman una base de~$\mathbb{R}^3$?

\n Determinar si la familia de
vectores $\{(1,1+2i)^T,(i,-2+i)^T\}$ es una base del espacio vectorial complejo
 $\mathbb{C}^2$.


%%%%%%%%%%%%%%%%%%%%%%%%%%%%%%%%%%%%%%%%%%%%%%%%%%%%%%%%%%%%%%%%%%%%%%%%%%%%%%%%%%%%%

\newpage %                   !!!!!!!!!!!!!!!!!!!!!!!!!!!!!!!!!!!!!!!!!!!

%%%%%%%%%%%%%%%%%%%%%%%%%%%%%%%%%%%%%%%%%%%%%%%%%%%%%%%%%%%%%%%%%%%%%%%%%%%%%%%%%%%%%


\n \A
Se considera el espacio vectorial real $\rpol{3}$.\\[-1.5\baselineskip]
\begin{enumerate}\itemsep=-2pt
\item ?`Es $W_1=\{ p(x)\in\rpol{3}\colon (x-1) \mbox{ divide }
p(x)\}$ subespacio de $\rpol{3}$?

\item ?`Es $W_2=\{ p(x)\in\rpol{3}\colon p(2)=0\}$ subespacio de
$\rpol{3}$?

\item Escribir, si es posible, los polinomios $1,x,x^2,x^3$ como
combinaci\'on lineal del sistema  de vectores formado por
$p(x)=1+x+x^2+x^3$ y sus tres primeras derivadas $p'(x), p''(x)$ y
$p'''(x)$.
\end{enumerate}


\n \A
Se consideran los subconjuntos $V$ y $W$ de $\C^4$ donde
$$
    V=\big\{(z_1,z_2,z_3,z_4)^T \in \C^4\colon\quad z_1=-i z_2+(1+i)z_3+z_4\big\}
$$
y las ecuaciones param�tricas de $W$ son
\vskip-2mm
$$
\left.
    \begin{array}{lll}
   z_ 1 &   =&\mu+i\beta \\
   z_ 2 &   =&\lambda+i\mu \\
      z_3  &   =&i\lambda \\
         z_4 &   =&\lambda+i\beta \\
    \end{array}
\right\}\, ,\hspace{1cm} \lambda,\,\mu,\,\beta\in\C.
$$
\begin{enumerate}

\item Comprobar si $V$ y $W$ son subespacios de $\C^4$.

\item Comprobar si es cierto que $W\subset V$.

\item Responder a las mismas preguntas, si en las ecuaciones param�tricas de
$W$ s�lo se admiten $\lambda,\,\mu,\,\beta$ reales.

\end{enumerate}


\n \A
El vector $v$ del espacio vectorial $\R^2$ tiene
coordenadas $(1,2)$ respecto de la base $B=\{v_1,v_2\}$.
Hallar las coordenadas de $v$ en la base $B'=\{u_1,u_2\}$ en $\R^2$,
sabiendo que
$$
    v_1=3u_1+2u_2;
    \qquad
    v_2=4u_1-v_2\,.
$$

\n \A
Estudiar cu\'ales de las siguientes aplicaciones son
lineales entre los espacios vectoriales dados:\\[-1\baselineskip]
\begin{enumerate}
\renewcommand{\labelenumi}{
%    \selectlanguage{english}
    \makebox[30pt][r]{\roman{enumi})}
    }
    \itemsep=-2pt
\item  $f_1\colon M_2(\R)\to M_2(\R)$ dada por $f_1(A)=AB$, donde
$B=\left(\begin{array}{cc}1 &- 1\\1 &-1\end{array}\right)$.

\item $f_2\colon M_2(\R)\to M_2(\R)$ dada por $f_2(A)=A+B$, con
$B\in M_2(\R)$ fija.

\item  $f_3\colon M_2(\R)\to M_2(\R)$ dada por $f_3(A)=AB-BA$,
donde $B=\begin{pmatrix}
    2  &  1  \\
    0  &  2
\end{pmatrix}$.

\item  $f_4\colon M_2(\C)\to V$ dada por $f_4(A)=A+A^T$, donde
$V=\{A\in M_2(\C)\colon A=A^T\}$.

\item $f_5\colon M_2(\R)\to W$ dada por $f_5(A)=AA^T$, donde
$W=\{A\in M_2(\R)\colon A=A^T\}$.

\item $f_6\colon \rpol{n}\to \rpol{n}$ dada por $f_6(p(x))=
p(x+1)$.

\item $f_7\colon \rpol{n}\to \rpol{n}$ dada por $f_7(p(x))=
p(x)+1$.

\item  $f_ 8\colon \R^3\to \R^3$ dada por $f_{8}(v)= \lambda_0v$
con $\lambda_0\in \R$ fijo.

\item  $f_ 9\colon \R^3\to \R^3$ dada por $f_{9}(v)= v_0-v $ con
$v_0\in \R^3$ fijo.

\item  $f_ {10}\colon \R^3\to \R^3$ dada por
$f_{10}(x_1,x_2,x_3)=(x_1,1,x_3)^T $.

\item  $f_{11} \colon \R^3\to \R^3$ dada por
$f_{11}(x_1,x_2,x_3)=(x_1^2-x_2^2,0,0)^T$.

\item  $f_ {12}\colon \R^3\to \R^3$ dada por
$f_{12}(x_1,x_2,x_3)=(x_1+x_2+x_3,x_1+x_2,x_3)^T $.

%\item  $f_ {13}\colon \R^3\to \R^2$ dada por $f_{13}(x_1,x_2,x_3)=
%(x_1-x_2,x_1+x_3)$.
\end{enumerate}











\end{document}
